% Part: many-valued-logic
% Chapter: infinite-valued-logics
% Section: goedel

\documentclass[../../../include/open-logic-section]{subfiles}

\begin{document}

\olfileid[tr]{mvl}{inf}{god}

\olsection{G\"odel Mantıkları}

\begin{editorial}
  Bu, sonsuz değerli G\"odel mantığına ilişkin kısa bir bölüm taslağıdır.
\end{editorial}

\begin{defn}\ollabel{def:goedel} Sonsuz değerli G\"odel
mantığı~$\LogGod[\infty]$ şu matrisle tanımlanır:
\begin{enumerate}
  \item $\lfalse$, $\lnot$, $\land$, $\lor$, $\lif$ bağlaçlarını içeren
  standart önerme dili $\Lang L_0$.
  \item Doğruluk değerleri kümesi $V_\infty$.
  \item Tek belirlenmiş değer $1$'dir; yani $V^+ = \{1\}$.
  \item Doğruluk fonksiyonları aşağıdaki fonksiyonlarla verilir:
  \begin{align*}
    \tf{\lfalse} & = 0\\
    \tf{\lnot}[\LogGod](x) & = \begin{cases}
      1 & \text{$x = 0$ ise}\\
      0 & \text{aksi durumda}
    \end{cases}\\
    \tf{\land}[\LogGod](x,y) & = \min(x,y)\\
    \tf{\lor}[\LogGod](x,y) & = \max(x,y)\\
    \tf{\lif}[\LogGod](x,y) & = \begin{cases}
      1 & \text{$x \le y$ ise}\\
      y & \text{aksi durumda.}
    \end{cases}
    \end{align*}
\end{enumerate}
$m$-değerli G\"odel mantığı da aynı biçimde, ancak $V = V_m$ alınarak
tanımlanır.
\end{defn}

\begin{prop}
  \olref[thr][god]{defn:goedel} ile tanımlanan $\LogGod[3]$ mantığı,
  \olref{def:goedel} ile tanımlanan $\LogGod[3]$ ile aynıdır.
\end{prop}

\begin{proof}
  Bu, \olref[thr][god]{defn:goedel} içindeki bağlaç doğruluk tabloları ile
  \olref{def:goedel} içindeki denklemlerin belirlediği doğruluk tabloları
  karşılaştırılarak görülebilir:
  \begin{center}
    \begin{tabular}{c|c} 
      $\tf{\lnot}[\LogGod[3]]$ & \\ 
      \hline  
      $1$ & $0$ \\ 
      $1/2$ & $0$ \\
      $0$ & $1$ 
    \end{tabular}
    \quad
    \begin{tabular}{c|ccc} 
      $\tf{\land}[\LogGod]$ & $1$ & $1/2$ & $0$ \\ 
      \hline 
      $1$ & $1$ & $1/2$ & $0$ \\ 
      $1/2$ & $1/2$ & $1/2$ & $0$\\ 
      $0$ & $0$ & $0$ & $0$ 
    \end{tabular}
    \\[2ex]
    \begin{tabular}{c|ccc} 
      $\tf{\lor}[\LogGod]$ & $1$ & $1/2$ & $0$ \\ 
      \hline 
      $1$ & $1$ & $1$ & $1$ \\ 
      $1/2$ & $1$ & $1/2$ & $1/2$ \\
      $0$ & $1$ & $1/2$ & $0$ 
    \end{tabular}
    \quad
    \begin{tabular}{c|ccc} 
      $\tf{\lif}[\LogGod]$ & $1$ & $1/2$ & $0$ \\ 
      \hline 
      $1$ & $1$ & $1/2$ & $0$ \\ 
      $1/2$ & $1$ & $1$ & $0$  \\ 
      $0$ & $1$ & $1$ & $1$ 
    \end{tabular}
  \end{center} 
\end{proof}

\begin{prop}\ollabel{prop:god-infty-m}
  $\Gamma \Entails[\LogGod[\infty]] !B$ ise bütün $m \ge 2$ değerleri için
  $\Gamma \Entails[\LogGod[m]] !B$.
\end{prop}

\begin{proof}
  Alıştırma.
\end{proof}

\begin{prob}
  \olref[mvl][inf][god]{prop:god-infty-m} önermesini ispatlayın.
\end{prob}

Aslında bunun tersi de geçerlidir.

$\LogGod[3]$ gibi $\LogGod[\infty]$ da sezgisel olarak geçerli bütün
!!{formula}s{}i totoloji olarak içerir; aynı totoloji-olmama örnekleri
$\LogGod[\infty]$'da da totoloji değildir:
\begin{align*}
  & p \lor \lnot p && (p \lif q) \lif (\lnot p \lor q) \\
  & \lnot\lnot p \lif p && \lnot(\lnot p \land \lnot q) \lif (p \lor q) \\
  & ((p \lif q) \lif p) \lif p && \lnot(p \lif q) \lif (p \land \lnot q)
\end{align*}
Sezgisel olarak geçersiz bir !!{formula} olmakla birlikte $\LogGod[3]$'ün
bir totolojisi olan $(p \lif q) \lor (q \lif p)$ örneği
$\LogGod[\infty]$'da da bir totolojidir. Gerçekte $\LogGod[\infty]$,
kendisine
$(!A \lif !B) \lor (!B \lif !A)$ şeması eklenmiş sezgisel mantık olarak
nitelenebilir. Bunu Michael Dummett göstermiştir; dolayısıyla
$\LogGod[\infty]$'a çoğu kez G\"odel--Dummett mantığı~$\Log{LC}$ denir.

\begin{prob}
  $(p \lif q) \lor (q \lif p)$'nin $\LogGod[\infty]$'ın bir totolojisi
  olduğunu gösterin.
\end{prob}

\begin{prob}
  $\LogGod[3]$'ün bir totolojisi olan
  $(p \lif q) \lor (q \lif r) \lor (r \lif s)$'nin
  $\LogGod[\infty]$'ın bir totolojisi olmadığını gösterin.
\end{prob}

\end{document}
